\newpage
\begin{center}
    {\large\bfseries LEMBAR WAWANCARA PRA-PENELITIAN}\\
    \vspace{0.2cm}
    {\bfseries (Untuk Guru Matematika)}
\end{center}

\vspace{0.5cm}

\begin{table}[H]
    \begin{tabular}{@{}ll@{}}
        Nama Guru      & : Wibowo Ramadhiyanto, S.Pd. \\
        Nama Sekolah   & : SMP Muhammadiyah 9 Yogyakarta \\
        Hari/Tanggal   & : Kamis, 30 April 2026 \\
    \end{tabular}
\end{table}

\vspace{0.3cm}

\noindent \textbf{Petunjuk}: Jawablah pertanyaan berikut sesuai dengan pengalaman Bapak/Ibu dalam pembelajaran matematika di kelas VIII.

\vspace{0.3cm}

\begin{longtable}{|c|p{8cm}|p{6cm}|}
    \hline
    \textbf{No.} & \textbf{Pertanyaan} & \textbf{Jawaban} \\
    \hline
    1 & Apakah Bapak/Ibu pernah menggunakan media pembelajaran berbasis teknologi (komputer, proyektor, simulasi 3D, AR, dll.) dalam pembelajaran matematika? Jika pernah, media apa saja yang digunakan? & \\
    \hline
    2 & Menurut Bapak/Ibu, materi matematika apa yang paling sulit dipahami oleh siswa kelas VIII? &  \\
    \hline
    3 & Khusus untuk materi bangun ruang sisi datar (kubus, balok, prisma, limas), kesulitan apa yang paling sering dialami siswa? & \\
    \hline
    4 & Selama ini, media pembelajaran apa yang biasa Bapak/Ibu gunakan untuk mengajarkan materi bangun ruang sisi datar? &  \\
    \hline
    5 & Apakah media yang digunakan selama ini sudah membantu siswa dalam memvisualisasikan bentuk tiga dimensi bangun ruang? & \\
    \hline
    6 & Apakah Bapak/Ibu pernah menggunakan atau mengembangkan media pembelajaran berbasis simulasi 3D atau \textit{Augmented Reality} (AR)? &  \\
    \hline
    7 & Menurut Bapak/Ibu, apakah pengembangan media pembelajaran berbasis simulasi 3D diperlukan untuk materi bangun ruang sisi datar? Mengapa? &  \\
    \hline
    8 & Kurikulum apa yang sedang diterapkan di sekolah untuk kelas VIII? & \\
    \hline
\end{longtable}

\vspace{0.5cm}

\begin{flushright}
    \begin{minipage}{0.48\textwidth}
        \raggedright
        Yogyakarta, 30 April 2026 \\
        Pewawancara,\\[1.5cm]
        \underline{Toriq Afanudin} \\
        NIM 1900006105
    \end{minipage}
    \hfill
    \begin{minipage}{0.48\textwidth}
        \centering
        Guru Matematika,\\[1.5cm]
        \underline{\hspace{6cm}}\\
    \end{minipage}
\end{flushright}